\section{Conclusions}\label{sec:conclusion}

Random parameters models have changed how unobserved heterogeneity is handled in crash frequency analysis, but the choice of which parameters are random is still made by sequential testing or by searching for a single best specification. This paper replaces that step with a Bayesian model in which each candidate covariate is excluded, fixed or random, and in which the posterior probabilities of these states are estimated jointly with the parameters. The identification results and the simulations explain why random parameters should be interpreted with care in cross-sectional and panel data alike. The sampler makes the approach practical by integrating out the site-specific draws when deciding whether a parameter is random.

In panel data the method recovered the full specification far more often than sequential testing or an information criterion search, and its intervals reflected the uncertainty about which parameters are random. Three practical recommendations follow for agencies developing safety performance functions. First, include squared covariates as fixed-only candidates, because omitted curvature otherwise passes for heterogeneity in cross-sectional and panel data alike. Second, use the false discovery rate rule when the claim of interest is that heterogeneity is present. Third, let the dispersion depend on the candidate covariates and select them in the same model, because a dispersion that varies with a covariate otherwise produces a random parameter on that covariate. In panel data this separates the two explanations. In cross-sectional data it can only divide the posterior between them, and a random parameter selected there should be read as evidence of heterogeneity in either the effect or the dispersion.

The application to 17{,}139 British motorway and A-road links shows these mechanisms in real data. The sequential procedure declared 13 of 15 covariates random. The proposed method declared nine, and modeling the mean curvature and the dispersion changed which effects it found random: two random parameters under a linear mean were explained by curvature, and part of the heterogeneity in another by a covariate-dependent dispersion. The resulting model predicted the next year's collisions better than every alternative. With data of this size almost every candidate is included and most random parameters have probabilities near zero or one, so the value of the method lies less in quantifying uncertainty about a single random parameter than in testing heterogeneity against the alternative explanations that a conventional specification leaves out.
